% !PRINTED 218
\weberpage{218}

$\Theta$ is a so-called decomposable form of degree $\frac{1}{2}n(n-1)$ in the variables $t$; for by (7) it can be decomposed entirely into linear factors, which, to be sure, are not rational in the $a$.

We have already pointed out in \S.\,52 that in the Tschirnhausen transformation $x$ is also expressible rationally in terms of $y$; and the same therefore holds for every rational function of $x$ as well.

If we consider any such function $\varphi(x)$, which may also contain the coefficients $a$ and $t$, but is always assumed to be integral and homogeneous in both kinds of variables, then we may set

\weq{9}{\varphi(x) = C_0 + C_1y + C_2y^2 + \cdots + C_{n-1}y^{n-1},}

in which the $C_0, C_1, \ldots C_{n-1}$ are expressible rationally in $a$ and in $t$.

If we write down equation (9) for $x=x_1,x_2\ldots x_n$, $y=y_1,y_2,\ldots y_n$, we obtain for the determination of the $C$ a system of linear equations whose determinant

\weq{10}{\begin{vmatrix} 1, & y_1, & y_1^2 & \ldots & y_1^{n-1} \\ 1, & y_2, & y_2^2 & \ldots & y_2^{n-1} \\ \cdot & \cdot & \cdot & \cdot & \cdot \\ 1, & y_n, & y_n^2 & \ldots & y_n^{n-1} \end{vmatrix},}

is therefore equal to $\Theta\sqrt{D}$. The numerator of the expression for $C_\nu$ arises from the determinant (10) by replacing the elements of the $(\nu+1)$th column by $\varphi(x_1), \varphi(x_2)\ldots\varphi(x_n)$, and consequently the numerator, when everything is expressed in terms of the $x_i$, has the factor $\sqrt{D}$ (because it vanishes when two of the $x_i$ become equal to each other).

We may therefore set

\[ \Theta C_\nu = Q_\nu, \]

in which $Q_\nu$ is an integral rational function of the $t$ and the $a$, which can at most still contain a power of $a_0$ in the denominator. We then obtain

\weq{11}{\Theta\varphi(x) = Q_0 + Q_1y + Q_2y^2 + \cdots + Q_{n-1}y^{n-1}.}

The coefficients $Q_\nu$, however, are not invariants, and their computation is in most cases difficult.
